\section{What \texttt{main.tex} does}

Philosophy aside, the spine of the build is almost insultingly small: one file lists the whole book in the order a reader meets it.

The file \verb|main.tex| is the book's driver. It chooses the document class, loads the shared layout and music helpers, and then pulls in each piece of front matter, each chapter, and the back matter with \verb|\input{...}| in reading order. Keep this file mostly structural: headings, \verb|\input| lines, and the occasional one-line note—not long finished prose.

\textbf{Where the music lives.} Engraved examples are not typed into \verb|main.tex|. Each musical fragment normally lives in its own \verb|.ly| file under the chapter tree (inline, excerpt, or full-page areas). The surrounding wrapper should usually live directly in the real chapter \verb|.tex| file: five or so lines with \verb|\lilypondfile{./path/to/score.ly}| inside the book's music helpers (\verb|\beginBookExample| and friends), not an extra throwaway \verb|.tex| wrapper just to hold one import. The \verb|lilypond-book| step turns LilyPond into graphics and LaTeX glue; \verb|make check-ly| can still run plain LilyPond on the same \verb|.ly| files for fast error checking.

When you grow the book, add a chapter folder, keep prose in section-sized \verb|.tex| files, keep each musical example in a \verb|.ly| you import, and add one \verb|\input| per prose file to \verb|main.tex| in the order readers should meet them—using ordinary \verb|\chapter| and \verb|\section| headings as you already do in the template.
